% ---------------------------------------------------------------------
% Vorlage f�r Dokumente, die auf report oder book basieren
% Voraussetzung: spezifische Variable wie \imagePath
% oder \theSemester sind bereits definiert.
% ---------------------------------------------------------------------
\typeout{Einstellung f�r Texte auf report/book-Basis}
\typeout{     (C) Manfred Brill}
\typeout{     Version 1.1 Juli 2018}
\usepackage{mb}
\usepackage{mbmath}
\usepackage{textcomp}
% Array-Paket f�r mehr Kontrolle der Tabellen
\usepackage{array}
% ifthen f�r Ein- und Ausblenden der L�sungen.
\usepackage{ifthen}
% Paket f�r Postscript Pi-Fonts
\usepackage{pifont}
% PSTricks
\usepackage{pstricks}
% Paket f�r das einfache Umdefinieren von Listen
\usepackage[shortlabels]{enumitem}
% Alphabetischer Index
\usepackage{imakeidx}
\makeindex[title=Index,columns=2,options=-s myGerman, intoc=true]
% Kopf- und Fusszeilen mit KoMaScript
\usepackage[automark, headsepline]{scrlayer-scrpage}
%
\pagestyle{scrheadings}
\clearscrheadfoot
% Kopfzeile definieren
\ihead{\headmark}
\ohead[]{\pagemark}
\chead{}
\pagestyle{scrheadings}
% Zentral definierte Variablen
% Abk�rzungen, Namen, ...
% -------------------------------------------------------------------
% variablen.tex
% Variablen f�r Werkzeuge, Begriffe, Firmen
% Stand: Juni 2017
% Liegt zentral in texmf-local/tex/latex/MB
% -------------------------------------------------------------------
% Programmiersprachen
\newcommand{\java}{\texttt{Java}}
\newcommand{\cpp}{\texttt{C++}}
\newcommand{\cl}{\texttt{C}}
\newcommand{\csharp}{\texttt{C\#}}
\newcommand{\php}{\texttt{PHP}}
\newcommand{\py}{\texttt{Python}}
\newcommand{\perl}{\texttt{Perl}}
\newcommand{\fc}{\texttt{Fortran}}
\newcommand{\obc}{\texttt{Objective C}}
\newcommand{\js}{\texttt{JavaScript}}
\newcommand{\omp}{\texttt{OpenMP}}
\newcommand{\mpi}{\texttt{MPI}}
\newcommand{\ocl}{\texttt{OpenCL}}
\newcommand{\cuda}{\texttt{CUDA}}
% Datei-Formate
\newcommand{\html}{\texttt{HTML}}
\newcommand{\xml}{\texttt{XML}}
\newcommand{\json}{\texttt{JSON}}
\newcommand{\rtf}{\texttt{RTF}}
\newcommand{\pdf}{\texttt{PDF}}
%
\newcommand{\sccs}{\texttt{SCCS}}
\newcommand{\rcs}{\texttt{RCS}}
\newcommand{\cvs}{\texttt{CVS}}
\newcommand{\svn}{\texttt{SVN}}
\newcommand{\subversion}{\texttt{Subversion}}
\newcommand{\git}{\texttt{Git}}
\newcommand{\github}{\texttt{GitHub}}
\newcommand{\githubDesktop}{\texttt{GitHub Desktop}}
\newcommand{\gitkraken}{\texttt{GitKraken}}
\newcommand{\gittortoise}{\texttt{TortoiseGit}}
\newcommand{\gittree}{\texttt{Sourcetree}}
\newcommand{\bitbucket}{\texttt{BitBucket}}
\newcommand{\mercurial}{\texttt{Mercurial}}
\newcommand{\pf}{\texttt{Perforce}}
\newcommand{\pfv}{\texttt{P4V}}
\newcommand{\pfc}{\texttt{p4}}
\newcommand{\pfmerge}{\texttt{P4Merge}}
\newcommand{\pfsand}{\texttt{P4.Sandbox}}

\newcommand{\vi}{\texttt{vi}}
\newcommand{\emacs}{\texttt{Emacs}}
%
\newcommand{\javadoc}{\texttt{Javadoc}}
\newcommand{\xmldoc}{\texttt{XMLdoc}}
\newcommand{\sandcastle}{\texttt{Sandcastle}}
\newcommand{\doxy}{\texttt{Doxygen}}
\newcommand{\doxyVersion}{\texttt{$1.8.7$}}
\newcommand{\doxywizard}{\texttt{Doxywizard}}
\newcommand{\graphviz}{\texttt{Graphviz}}
%
\newcommand{\logj}{\texttt{Log4j}}
\newcommand{\slfj}{\texttt{SLF4J}}
\newcommand{\logback}{\texttt{LOGBack}}
\newcommand{\jul}{\texttt{Java Logging API}}
\newcommand{\commonsLog}{\texttt{Apache Commons LoggingJ}}
\newcommand{\boostlog}{\texttt{Boost.Log}}
\newcommand{\glog}{\texttt{Google Logging Library}}
%
\newcommand{\junit}{\texttt{JUnit}}
\newcommand{\xunit}{\texttt{xUnit}}
\newcommand{\nunit}{\texttt{NUnit}}
\newcommand{\xunitnet}{\texttt{xUnit.net}}
\newcommand{\ctest}{\texttt{CTest}}
%
\newcommand{\ssh}{\texttt{ssh}}
\newcommand{\sftp}{\texttt{sftp}}
\newcommand{\make}{\texttt{make}}
\newcommand{\ant}{\texttt{Ant}}
\newcommand{\maven}{\texttt{Maven}}
\newcommand{\gradle}{\texttt{Gradle}}
\newcommand{\msbuild}{\texttt{MSBuild}}
\newcommand{\mstest}{\texttt{MSTest}}
\newcommand{\cmake}{\texttt{CMake}}
\newcommand{\jenkins}{\texttt{Jenkins}}
\newcommand{\ccontrol}{\texttt{cruisecontrol}}
\newcommand{\cdash}{\texttt{CDash}}
%
\newcommand{\eclipse}{\texttt{Eclipse}}
\newcommand{\vs}{\texttt{Visual Studio}}
\newcommand{\mono}{\texttt{Mono}}
\newcommand{\monodev}{\texttt{MonoDevelop}}
%
\newcommand{\openO}{\texttt{OpenOffice}}
\newcommand{\Qt}{\texttt{Qt}}
\newcommand{\firefox}{\texttt{Mozilla Firefox}}
\newcommand{\kde}{\texttt{KDE}}
%
\newcommand{\ms}{\texttt{Microsoft}}
\newcommand{\net}{\texttt{.NET}}
\newcommand{\sun}{\texttt{SUN}}
\newcommand{\ora}{\texttt{Oracle}}
\newcommand{\apache}{\texttt{Apache}}
\newcommand{\windows}{\texttt{Windows}}
\newcommand{\cyg}{\texttt{Cygwin}}
\newcommand{\linux}{\texttt{Linux}}
\newcommand{\unix}{\texttt{Unix}}
\newcommand{\osx}{\texttt{MacOS X}}
\newcommand{\android}{\texttt{Android}}
% VR Software
\newcommand{\unity}{\texttt{Unity}}
\newcommand{\verUnity}{\texttt{$2018.2$}}
\newcommand{\unreal}{\texttt{Unreal}}
\newcommand{\verUnreal}{\texttt{4}}
\newcommand{\juggler}{\texttt{VRJuggler}}
\newcommand{\cavelib}{\texttt{CAVELib}}
\newcommand{\gadgeteer}{\texttt{Gadgeteer}}
\newcommand{\vrpn}{\texttt{VRPN}}
\newcommand{\mvr}{\texttt{MiddleVR}}
\newcommand{\verMvr}{\texttt{$1.7.0.7$}}
% Computergrafik und Werkzeuge
\newcommand{\direct}{\texttt{Direct3D}}
\newcommand{\gl}{\texttt{OpenGL}}
\newcommand{\para}{\texttt{ParaView}}
\newcommand{\vtk}{\texttt{VTK}}
% R und anderes zur Datenanalyse
% \R ist schon definiert (Mathematik!)
\newcommand{\Rsoft}{\texttt{R}}
\newcommand{\Rstudio}{\texttt{RStudio}}
\newcommand{\verR}{\texttt{$3.3.3$}}
\newcommand{\mondrian}{\texttt{Mondrian}}
\newcommand{\tableau}{\texttt{Tableau}}
% LaTeX
\newcommand{\jabref}{\texttt{JabRef}}
\newcommand{\bibtex}{\textsc{Bib}\negthinspace\TeX}

% Farben
% ---------------------------------------------------------------------------------------
% Farben, f�r Folien und andere Dinge.
% Version vom Juli 2015 mit neuer Fachbereichsfarbe
%
% Letzte �nderung: 1.9.2015
%
% Um diese Datei zu verwenden muss Sie in texmf-local/MB kopiert werden
% und TeX muss aktualisiert werden!
% ---------------------------------------------------------------------------------------
\definecolor{light}{gray}{.75}
\definecolor{LightGray}{rgb}{0.24,0.24,0.24}
%   Farben rot gr�n blau f�r Koordinatenachsen
\definecolor{wred}{rgb}{0.8, 0.0, 0.0}
\definecolor{wgreen}{rgb}{0.0, 0.8, 0.0}
\definecolor{wblue}{rgb}{0.0, 0.0, 0.8}
\definecolor{wyellow}{rgb}{1.0, 1.0, 0.0}
\definecolor{wcyan}{rgb}{0.0, 1.0, 1.0}
% Fachbereichsfarbe neu
\definecolor{colorDepartment}{rgb}{0.16, 0.71, 0.86}
% Listings
\definecolor{lstback}{gray}{0.85}
%

% Koordinatensysteme und andere Makros
% -----------------------------------------------------------------------------
%  Save-Boxes f�r Folien, in denen Koordinatensysteme und Karo-Muster
%  verwendet werden sollen.
%
%  Autor:   MB
%  History: 19.02.2009 Erstellt aus leeren Folien "Transformationen"
%           06.11.2012 Polarplots hinzugef�gt
%           26.01.2016 in texmf-local kopiert.
% -----------------------------------------------------------------------------
\newsavebox{\gridLines}
\savebox{\gridLines}{%
\begin{picture}(8.5,5.5)(-4.25,-2.2)
%\thicklines
\color[gray]{0.8}
\multiput(-4, -2)(0, 0.5){6}{\line(1,0){8.0}}
\multiput(-4, 0.5)(0, 0.5){6}{\line(1,0){8.0}}
\multiput(-4, -2)(0.5, 0){17}{\line(0,1){5.0}}
\color[gray]{0}
\end{picture}}% Ende \gridLines
%%%%%%%%%%%%%%%%%%%%%%%%%%%%%%%%%%%%%%%%%%%%%%%%%%%%%%%%%%%%%%%%%%%
\newsavebox{\centeredCS}
\savebox{\centeredCS}{%
\begin{picture}(8.5,5.5)(-4.25,-2.2)
%\thicklines
\color[gray]{0.8}
\multiput(-4, -2)(0, 0.5){6}{\line(1,0){8.0}}
\multiput(-4, 0.5)(0, 0.5){6}{\line(1,0){8.0}}
\multiput(-4, -2)(0.5, 0){8}{\line(0,1){5.0}}
\multiput(0.5, -2)(0.5, 0){8}{\line(0,1){5.0}}
\color[gray]{0}
\put(0,-2.2){\vector(0,1){5.5}}
\put(-4.25,0){\vector(1,0){8.5}}
\end{picture}}% Ende \centeredCS
%%%%%%%%%%%%%%%%%%%%%%%%%%%%%%%%%%%%%%%%%%%%%%%%%%%%%%%%%%%%%%%%%%%
\newsavebox{\positiveYCS}
\savebox{\positiveYCS}{%
\begin{picture}(8.5,4.5)(-4.25,-0.2)
%\thicklines
\color[gray]{0.8}
\multiput(-4, 0.5)(0, 0.5){8}{\line(1,0){8.0}}
\multiput(-4, 0.0)(0.5, 0){8}{\line(0,1){4.0}}
\multiput(0.5, 0.0)(0.5, 0){8}{\line(0,1){4.0}}
\color[gray]{0}
\put(0,-0.2){\vector(0,1){4.5}}
\put(-4.25,0){\vector(1,0){8.5}}
\end{picture}}% Ende \positiveYCS
%%%%%%%%%%%%%%%%%%%%%%%%%%%%%%%%%%%%%%%%%%%%%%%%%%%%%%%%%%%%%%%%%%%
\newsavebox{\firstquadrantCSsmall}
\savebox{\firstquadrantCSsmall}{%
\begin{picture}(8.6,5.5)(-0.2,-0.2)
\color[gray]{0.8}
\multiput(0, 0.25)(0, 0.25){20}{\line(1,0){8.0}}
\multiput(0.25, 0)(0.25, 0){32}{\line(0,1){5.0}}
\color[gray]{0}
\put(0,-0.2){\vector(0,1){5.5}}
\put(-0.2,0){\vector(1,0){8.6}}
\end{picture}}% Ende \firstquadrantCSsmall
%%%%%%%%%%%%%%%%%%%%%%%%%%%%%%%%%%%%%%%%%%%%%%%%%%%%%%%%%%%%%%%%%%%
\newsavebox{\firstquadrantCS}
\savebox{\firstquadrantCS}{%
\begin{picture}(8.6,5.5)(-0.2,-0.2)
\color[gray]{0.8}
\multiput(0, 0.5)(0, 0.5){10}{\line(1,0){8.0}}
\multiput(0.5, 0)(0.5, 0){16}{\line(0,1){5.0}}
\color[gray]{0}
\put(0,-0.2){\vector(0,1){5.5}}
\put(-0.2,0){\vector(1,0){8.6}}
\end{picture}}% Ende \firstquadrantCS
%%%%%%%%%%%%%%%%%%%%%%%%%%%%%%%%%%%%%%%%%%%%%%%%%%%%%%%%%%%%%%%%%%%
\newsavebox{\axis}
\savebox{\axis}{%
\begin{picture}(8.5,5.5)(-4.2,-2.75)
%\thicklines
\color[gray]{0}
\put(0,-2.75){\vector(0,1){5.5}}
\put(-4.25,0){\vector(1,0){8.5}}
\end{picture}}% Ende \axis
%%%%%%%%%%%%%%%%%%%%%%%%%%%%%%%%%%%%%%%%%%%%%%%%%%%%%%%%%%%%%%%%%%%
\newsavebox{\axisWithNumbers}
\savebox{\axisWithNumbers}{%
\begin{picture}(8.5,5.5)(-4.2,-2.75)
%\thicklines
\color[gray]{0}
\put(0,-2.75){\vector(0,1){5.5}}
\put(-4.25,0){\vector(1,0){8.5}}
% Beschriftung
\multiput(1, -0.1)(1,0){3}{\line(0,1){0.2}}
\multiput(-3, -0.1)(1,0){3}{\line(0,1){0.2}}
\multiput(-0.1, 1)(0,1){2}{\line(1,0){0.2}}
\multiput(-0.1, -2)(0,1){2}{\line(1,0){0.2}}
%
\small
\put(0.95, -0.4){$1$}
\put(-0.4, 0.95){$1$}
\end{picture}}% Ende \axisWithNumbers
%%%%%%%%%%%%%%%%%%%%%%%%%%%%%%%%%%%%%%%%%%%%%%%%%%%%%%%%%%%%%%%%%%%
\newsavebox{\shortAxis}
\savebox{\shortAxis}{%
\begin{picture}(3.4,3.4)(-1.7,-1.7)
%\thicklines
\color[gray]{0}
\put(0,-1.7){\vector(0,1){3.4}}
\put(-1.7,0){\vector(1,0){3.4}}
\end{picture}}% Ende \shortAxis
%%%%%%%%%%%%%%%%%%%%%%%%%%%%%%%%%%%%%%%%%%%%%%%%%%%%%%%%%%%%%%%%%%%
\newsavebox{\shortAxisWithNumbers}
\savebox{\shortAxisWithNumbers}{%
\begin{picture}(3.4,3.4)(-1.7,-1.7)
%\thicklines
\color[gray]{0}
\put(0,-1.7){\vector(0,1){3.4}}
\put(-1.7,0){\vector(1,0){3.4}}
% Beschriftung
\put(1, -0.1){\line(0,1){0.2}}
\put(-1, -0.1){\line(0,1){0.2}}
\put(-0.1, 1){\line(1,0){0.2}}
\put(-0.1, -1){\line(1,0){0.2}}
%
\small
\put(0.95, -0.4){$1$}
\put(-0.4, 0.95){$1$}
\end{picture}}% Ende \shortAxisWithNumbers
%%%%%%%%%%%%%%%%%%%%%%%%%%%%%%%%%%%%%%%%%%%%%%%%%%%%%%%%%%%%%%%%%%%
\newsavebox{\positiveAxis}
\savebox{\positiveAxis}{%
\begin{picture}(5.5,5.5)(-0.2,-0.2)
\put(0,-0.2){\vector(0,1){5.5}}
\put(-0.2,0){\vector(1,0){5.5}}
\end{picture}}% Ende \positiveAxis
%%%%%%%%%%%%%%%%%%%%%%%%%%%%%%%%%%%%%%%%%%%%%%%%%%%%%%%%%%%%%%%%%%%
\newsavebox{\shortPositiveAxis}
\savebox{\shortPositiveAxis}{%
\begin{picture}(3.5,3.5)(-0.2,-0.2)
\put(0,-0.2){\vector(0,1){3.5}}
\put(-0.2,0){\vector(1,0){3.5}}
\end{picture}}% Ende \shortPositiveAxis
%%%%%%%%%%%%%%%%%%%%%%%%%%%%%%%%%%%%%%%%%%%%%%%%%%%%%%%%%%%%%%%%%%%
\newsavebox{\positiveAxisWithNumbers}
\savebox{\positiveAxisWithNumbers}{%
\begin{picture}(5.5,5.5)(-0.2,-0.2)
\put(0,-0.2){\vector(0,1){5.5}}
\put(-0.2,0){\vector(1,0){5.5}}
% Beschriftung
\multiput(1, -0.1)(1,0){4}{\line(0,1){0.2}}
\multiput(-0.1, 1)(0,1){4}{\line(1,0){0.2}}
%
\small
\put(0.95, -0.4){$1$}
\put(-0.4, 0.95){$1$}
\end{picture}}% Ende \positiveAxisWithNumbers
%%%%%%%%%%%%%%%%%%%%%%%%%%%%%%%%%%%%%%%%%%%%%%%%%%%%%%%%%%%%%%%%%%%
\newsavebox{\polarCircles}
\savebox{\polarCircles}{%
\begin{picture}(6.5,6.5)(-3.25,-3.25)
%\thicklines
\color[gray]{0.8}
% Argument bei circle ist der Durchmesser!
\put(0,0){\circle{2.0}}
\put(0,0){\circle{4.0}}
\put(0,0){\circle{6.0}}
\color[gray]{0}
\end{picture}}% Ende \polarCircles
%%%%%%%%%%%%%%%%%%%%%%%%%%%%%%%%%%%%%%%%%%%%%%%%%%%%%%%%%%%%%%%%%%%
\newsavebox{\polarCirclesFine}
\savebox{\polarCirclesFine}{%
\begin{picture}(6.5,6.5)(-3.25,-3.25)
%\thicklines
\color[gray]{0.8}
% Argument bei circle ist der Durchmesser!
\put(0,0){\circle{1.0}}
\put(0,0){\circle{2.0}}
\put(0,0){\circle{3.0}}
\put(0,0){\circle{4.0}}
\put(0,0){\circle{5.0}}
\put(0,0){\circle{6.0}}
\put(0,0){\circle{7.0}}
\color[gray]{0}
\end{picture}}% Ende \polarCirclesFine
%%%%%%%%%%%%%%%%%%%%%%%%%%%%%%%%%%%%%%%%%%%%%%%%%%%%%%%%%%%%%%%%%%%
\newsavebox{\polarAxis}
\savebox{\polarAxis}{%
\begin{picture}(6.5,6.5)(-3.25,-3.25)
\put(-2.12,-2.12){\line(1,1){4.24}}
\put(2.12, -2.12){\line(-1, 1){4.24}}
\thicklines % damit die "normalen" Striche auch gleich dick sind ...
\put(-3,0){\line(1,0){6.0}}
\put(0,-3){\line(0,1){6.0}}
\end{picture}}% Ende \polarAxis
%%%%%%%%%%%%%%%%%%%%%%%%%%%%%%%%%%%%%%%%%%%%%%%%%%%%%%%%%%%%%%%%%%%
\newsavebox{\polarAxisFine}
\savebox{\polarAxisFine}{%
\begin{picture}(6.5,6.5)(-3.25,-3.25)
\put(-2.47,-2.47){\line(1,1){4.94}}
\put(2.47, -2.47){\line(-1, 1){4.94}}
\thicklines % damit die "normalen" Striche auch gleich dick sind ...
\put(-3.5,0){\line(1,0){7.0}}
\put(0,-3.5){\line(0,1){7.0}}
\end{picture}}% Ende \polarAxisFine
%%%%%%%%%%%%%%%%%%%%%%%%%%%%%%%%%%%%%%%%%%%%%%%%%%%%%%%%%%%%%%%%%%%
\newsavebox{\polarPlot}
\savebox{\polarPlot}{%
\begin{picture}(6.5,6.5)(-3.25,-3.25)
\put(-3.25,-3.25){\usebox{\polarAxis}}
\put(-3.25,-3.25){\usebox{\polarCircles}}
\end{picture}}% Ende \polarPlot
%%%%%%%%%%%%%%%%%%%%%%%%%%%%%%%%%%%%%%%%%%%%%%%%%%%%%%%%%%%%%%%%%%% \newsavebox{\polarPlot}
\newsavebox{\polarPlotFine}
\savebox{\polarPlotFine}{%
\begin{picture}(6.5,6.5)(-3.25,-3.25)
\put(-3.25,-3.25){\usebox{\polarAxisFine}}
\put(-3.25,-3.25){\usebox{\polarCirclesFine}}
\end{picture}}% Ende \polarPlotFine
%%%%%%%%%%%%%%%%%%%%%%%%%%%%%%%%%%%%%%%%%%%%%%%%%%%%%%%%%%%%%%%%%%%

%
\raggedbottom
\setlength{\parskip}{2.0ex}
\setlength{\parindent}{0.0cm}
% Verhindert Schusterjungen und Hurenkinder
\clubpenalty = 10000
\widowpenalty = 10000
\displaywidowpenalty = 10000
%
\setcounter{secnumdepth}{2}
\setcounter{tocdepth}{2}
% Abst�nde zwischen Caption und Bild/Tabelle
\setlength\abovecaptionskip          {0.4em}
\setlength\belowcaptionskip          {0.2em}
% Anteil der Grafiken auf einer Seite sehr hoch setzen
% Damit wird verhindert, dass bei vielen Abbildungen im Text die
% Abbildungen ans Ende des PDF geschoben werden.
\renewcommand{\topfraction}{0.99}
\renewcommand{\textfraction}{0.0001}
% Literatur-Stil
\bibliographystyle{geralpha}
% Font und Fettdruck f�r Tabelle/Abbildung (neu, jetzt mit komascript!
% MB, 28/07/2017
\addtokomafont{caption}{\small}
\setkomafont{captionlabel}{\bfseries}
% Einstellung f�r Gliederungs�berschriften mit KoMaScript
\RedeclareSectionCommands[
  beforeskip=-.5\baselineskip,
  afterskip=.25\baselineskip
]{section,subsection,subsubsection}
%
% listings-Paket f�r Quelltexte
%++++++++++++++++++++++++++++++++++++++++++++++++++++++++++++++++
%   Hintergrundfarbe von Quellcode
\definecolor{codecolor}{rgb}{0.85,0.85,0.85}
\lstloadlanguages{Java}
\lstset{language=Java}
\definecolor{lstback}{gray}{0.85}
\lstset{backgroundcolor=\color{lstback}}
\lstset{extendedchars=true}
\lstset{showstringspaces = false}
\lstset{basicstyle = \ttfamily \small}
% Kommando f�r den Flattersatz bei nebeneinander liegenden Abbildungen
\newcommand*{\flatter}{\setlength{\rightskip}{0pt plus 2cm}}
% Anteil der Grafiken h�her auf jeder Seite!
\renewcommand{\textfraction}{0.001}
%
% Schritte in einer Aufz�hlung, daf�r einen Z�hler (schritt) und die Umgebung
% schritte definieren.
\newcounter{schritt}
\newenvironment{schritte}%
{\begin{list}%
{Schritt \arabic{schritt}:}%
{\usecounter{schritt}\settowidth{\labelwidth}{Schritt 1:}%
\setlength{\leftmargin}{\labelwidth}\addtolength\leftmargin{\labelsep}%
\parsep0.0ex\partopsep-0.3ex\itemsep2pt\topsep0.0ex}}{\end{list}}
%   Taste
\newcommand*{\taste}[1]{\small\textsf{#1}\normalsize}
%   gesch�tzte Namen (NICHT in chapter, section usw. verwenden!)
\newcommand*{\name}[1]{\textsl{#1}}
%   Fachbegriffe/Erkl�rung von Abk�rzungen
\newcommand*{\begriff}[1]{\index{#1}\textit{#1}}
%
\newcommand{\algorithmus}[2]{%
\vspace{4pt}\fboxsep 1mm \framebox[140mm]
{\parbox{135mm}{{\textbf{#1}}\vspace{2pt}#2}}\vspace{4pt}}
%
%   Dateinamen/Pfade
\newcommand*{\datei}[1]{\texttt{#1}\normalsize}
%   Tip (in einer Box)
\newcommand{\tip}[1]{%
\vspace{4pt}\fboxsep 1mm \framebox[140mm]
{\parbox{135mm}{{\textbf{Tipp}}:\\ #1}}\vspace{4pt}}
%
%   Reflektion (in einer Box)
\newcommand{\reflection}[1]{
\begin{quote}\fboxsep 3mm\framebox[140mm][c]{\parbox{130mm}{{\textbf Reflektion:\\}#1}}\end{quote}}
%
% Angabe im Begleittext, was die Studierenden lesen sollen (in einer Box)
\newcommand{\lesen}[1]{%
\begin{minipage}[c]{2.5cm}
\includegraphics[width=0.8\textwidth]{\generalImagePath/Misc/buchicon}%
\end{minipage}%
\begin{minipage}[t]{12cm}%
#1%
\end{minipage}}
%
%
%   Angabe im Begleittext, was die Studierenden lesen sollen (in einer Box)
%   Dieses Icon wird f�r Angaben verwendet, die nicht verpflichtend zu lesen
%   sind. Also "nice-to-have", "wenn sie noch Zeit haben".
%   Im Gegensatz zu lesen, was einen verpflichtenden Charakter hat.
\newcommand{\vertiefen}[1]{%
\begin{minipage}[b]{2.5cm}
\includegraphics[width=\textwidth]{\generalImagePath/Misc/reading}%
\end{minipage}%
\begin{minipage}[b]{12cm}%
#1%
\end{minipage}}
%
% Ein Counter f�r die Kontrollfragen. Wir z�hlen diesen Counter selbst hoch
% mit einem entsprechenden Kommando, das wir in \item verwenden.
\newcounter{kontrollCounter}[chapter]
\newcommand{\kontroll}[0]{
    \refstepcounter{kontrollCounter}
    % Kapitelnummer.Z�hler f�r das Item
    \arabic{chapter}.\arabic{kontrollCounter}
    % Das Label ist bis auf weiteres "kontrolle:Kapitelnummer:kontrollcounter
    \label{kontrolle:\arabic{chapter}:\arabic{kontrollCounter}}
}
% Definition der Kontrollfrage; die Nummerierung
% f�r das ganze erhalten wir mit Hilfe von \item{\kontroll}.
\newcommand{\kontrollfrage}[1]{%
\begin{minipage}[c]{1.85cm}
\huge{\ding{46}}%
\end{minipage}%
\begin{minipage}[t]{12cm}%
\begin{itemize}#1%
\end{itemize}%
\end{minipage}}
% Schalter f�r das ein- und ausblenden der Musterl�sungen
\newboolean{solutions}
%
% Dateinamen
\newcommand{\filename}[1]{%
\ifthenelse{\boolean{solutions}}{\framebox[50mm]{\parbox{40mm}{\textbf #1}}\vspace{6pt}}{}}
%
% Theorem-Umgebung f�r die �bungsaufgaben
% Wichtig: alle Attribute einstellen, dann die neue
% theorem-Umgebung mit newtheorem definieren!
% Oder, wie hier, durch das Einschlie�en in {}
%
{
% Zeilenumbruch bei Aufgaben-�berschrift
\theoremstyle{break}
% "Normaler" Font im Text
\theorembodyfont{\normalfont}
% Kapitelweise neu nummerieren
\newtheorem{auftitle}{Aufgabe}[chapter]
}
%
% Definition f�r die Kennzeichnung der �bungsaufgaben
%
\newcommand{\uebung}{%
\vspace*{11pt}%
\begin{tabular}{@{}p{2.25cm}@{}p{11.7cm}}%
\huge{\ding{45}}&\Large{\textbf{�bungsaufgaben}}%
\end{tabular}%
}
% Ein Counter f�r die �bungsaufgaben. Wir z�hlen diesen Counter selbst hoch
% mit einem entsprechenden Kommando, das wir in \item verwenden.
\newcounter{aufgabenCounter}[chapter]
\newcommand{\auf}[0]{
    \refstepcounter{aufgabenCounter}
    % Kapitelnummer.Z�hler f�r das Item
    \arabic{chapter}.\arabic{aufgabenCounter}
}
% Der Text der Aufgaben steht im Ordner ../Uebungsaufgaben/aufgaben/aufgabenstellungen.
% Die Struktur ist so wie in den Mathematikvorlesungen.

% Umgebungen f�r Satz, Definition, Beweis, ... . Wir orientieren uns am Mathematik-Buch,
% dort gab es diese Umgebungen auch schon. Im Grunde sind das einfach
% wieder theorem-Umgebungen.
{
\setlength\theorempreskipamount{5pt plus 3pt minus 1.5pt}
\setlength\theorempostskipamount{5pt plus 1pt minus 1pt}
% "Normaler" Font im Text
\theorembodyfont{\normalfont}
% Die Namen sind so gew�hlt, dass sie kompatibel zu Beamer sind; dann k�nnen
% wir den Text aus Folien kopieren und umgekehrt auch.
\newtheorem{Satz}{Satz}[chapter]
\newtheorem{Fakt}{Fakt}[chapter]
\newtheorem{Definition}{Definition}[chapter]
}
% Text mit Pfad, der vorher im Dokument definiert ist.
\newcommand{\aufgabentext}[1]{\renewcommand{\labelenumi}{\alph{enumi})}\auftitle\label{#1}\input{\exercisePath/aufgaben/aufgabenstellungen/#1}}
% Funktion f�r die L�sungs-Hinweise f�r eine Aufgabe. Der Text
% steht analog zu den �bungsaufgaben im Ordner aufgaben/loesungen.
\newcommand{\hinweistext}[1]{\renewcommand{\labelenumi}{\alph{enumi})}\input{\exercisePath/aufgaben/loesungen/#1}}
% Und jetzt die Funktionen, die wir im Text aufrufen
\newcommand{\aufgabe}[1]{\aufgabentext{#1}}
% L�sungshinweise im Anhang
\newcommand{\hinweis}[1]{\subsubsection*{Aufgabe \ref{#1}} \hinweistext{#1}}
%
% Marginalien
%
\newcommand{\randnotiz}[1]{\marginpar{\small{\textbf{#1}}}}
%
% Gabelschl�ssel als Marginalie und Hinweis auf Praxisbezug
\newcommand{\praxisbezug}[0]{\randnotiz{\includegraphics[width=1cm]{\generalImagePath/Misc/gabel}}}
%
% Smile Face
\newcommand{\smileface}[0]{\randnotiz{\includegraphics[width=1cm]{\generalImagePath/Misc/Face-smile}}}
%
% alert aus beamer-Folien zu emph machen
\newcommand*{\alert}[1]{\emph{#1}}
%
% Das obligatorische Argument ist der Titel des Dokuments.
% Als Default wird software-project.eps als Titelbild
% verwendet. Mit Hilfe eines optionalen Arguments kann
% ein anderes Bild verwendet werden!
% Beispiel: \titelseite[\imagePath/Misc/ameise]{Das Liebesleben der Ameisen}
%
\newcommand{\titelseite}[2][\titleImage]{%
%% Titelseite
\pagenumbering{roman}%
\thispagestyle{empty}%
\begin{titlepage}%
% Volle Zeilenbreite verwenden!
\centering%
\vspace*{0.5cm}%

\Huge{\theClass{}}\\\vspace*{0.5cm}%

\Huge{#2}%

\vspace*{3.0cm}%
\includegraphics[height=6cm]{#1}%

\vspace*{3.0cm}%
\Large{\theSchool{}}%

\Large{\theTeacher{}}%

\Large{\theSemester{}}%
\end{titlepage}}
